% Stage 8 restored fragment — 2.5.1 Source-to-field integral representations (archive lines 6202-6550)
\subsection{Source-to-field integral representations}
\label{sec:ch251}

Retarded source integrals are what method-of-moments codes discretize on antennas \cite{harrington2001}. Far-zone extraction is a limit on the same integral \cite{jackson1999}. Maxwell's retarded export to $\Gamma_R$ remains the governing constraint \cite{maxwell1873}.
The representation begins from retarded potentials driven by $\Jimp$ on compact $\Omega_s$, not from a plane-wave expansion \cite{maxwell1873}. The quantity being organized is the outgoing field pair $(\E,\H)$; angular labels attach only after $\Pi_{\mathrm{rad}}$ \cite{sommerfeld1949}. An alternative route through static potentials would omit radiation damping and misstate reactive storage \cite{jackson1999}. On a horn range, this step is what a MoM mesh implements before any pattern plot \cite{harrington2001}.

Before any volume quadrature is written, the range engineer already knows which observable the
campaign will report \cite{harrington2001}. A calibrated pattern probe on $\Gamma_R$ integrates
tangential $\E$ or $\H$ against a known polarization and yields a voltage proportional to a
directional amplitude---a discrete sample of the far-zone pattern carried by
$\mathcal{R}_\omega[\Jimp]$, not of the full pair $\mathcal{S}_\omega[\Jimp]$
\cite{stratton1941}. A shell wattmeter closes the same budget through
$\mathcal{P}_{\mathrm{rad}}$ on $\Gamma_R$ after $\Pi_{\mathrm{rad}}$ \cite{stratton1941}. Near-field
scanners on a hull enclosing $\Omega_s$ record $\mathcal{S}_\omega[\Jimp]$ and require a transform
before pattern meters apply \cite{hansen1988}. The volume integrals below are the constructive
route to those meter readings: they assemble $(\E,\H)$ from impressed $\Jimp$, after which export
selection maps the result to the pattern quantities the range records \cite{stratton1941}.
Radiation integrals assemble far-zone fields as superpositions of current moments weighted by
outgoing phase \cite{stratton1941}. Four laboratory settings motivate the
representations that follow. On an antenna range, probes on a large sphere $\Gamma_R$ sample
the field produced by a feed in the compact range; the measured voltage is a discretization of
a surface integral over $\Gamma_R$, not a direct reading of impressed $\Jimp$
\cite{harrington2001}. Horn illumination fixes $\Jimp$ on the mouth through a waveguide mode or
aperture integral; the pattern on the range wall appears only after those currents convolve
against the outgoing Green function \cite{collin2001}. Aperture radiation replaces the volume
support by equivalent currents on $\Gamma_{\mathrm{eq}}$ when the interior is inaccessible
\cite{stratton1941}. Near-field scanners reconstruct far-zone quantities from dense samples on
a hull enclosing $\Omega_s$---a transform that succeeds only when outgoing uniqueness and export
selection are respected \cite{hansen1988}.
The representation begins from retarded potentials driven by $\Jimp$ on compact $\Omega_s$,
organizes the outgoing pair $(\E,\H)$ before angular labels attach, and coincides with what
method-of-moments codes discretize before far-zone extraction
\cite{maxwell1873}. Static-potential routes omit radiation damping
and misstate reactive storage; they are not used here \cite{jackson1999}. Outgoing kernels,
regularization, and branch selection are presupposed from Sections~\ref{sec:ch23}--\ref{sec:ch24};
distance-regime asymptotics and flux meters follow in the subsections below.
Assumptions. Fix \(\omega>0\), compact \(\operatorname{supp}\Jimp=\Omega_s\Subset\Omega\), admissible
data Eq.~\eqref{eq:ch1.admissible-data-class}, charge continuity
Eq.~\eqref{eq:ch1.charge-continuity-harmonic}, and outgoing radiation condition
Eqs.~\eqref{eq:ch2.sommerfeld-outgoing-class}--\eqref{eq:ch2.outgoing-selector-def} on
exterior problems \cite{stratton1941}. Observation points \(\mathbf{r}\) lie outside
\(\Omega_s\) or on equivalent surfaces with trace matching declared below; coincidence
limits use Subsection~\ref{sec:ch232} \cite{harrington2001}. Kernels
\(\overline{\mathbf{G}}_{E,\omega}^{\mathrm{(out)}}\),
\(\overline{\mathbf{G}}_{H,\omega}^{\mathrm{(out)}}\), and scalar
\(G_\omega^{\mathrm{(out)}}\) denote outgoing branches fixed by
Eq.~\eqref{eq:ch2.outgoing-green-branch} \cite{stratton1941}.
The electric volume integral is the constructive form of the dyadic convolution route: each
source element \(\Jimp(\mathbf{r}')\,d^{3}r'\) contributes a dipole-like response that
propagates outward with retarded phase, and the total field is their coherent sum
\cite{stratton1941}. On a MoM mesh, patch currents discretize the same sum; on $\Gamma_R$, pattern
probes sample the far-zone limit after $\mathcal{R}_\omega$ has removed reactive storage that the
volume integral still carries at finite $r$ \cite{harrington2001}. In the symbols below,
$\mathbf{r}$ is the observation coordinate, $\mathbf{r}'$ the source coordinate integrated over
$\Omega_s$, $\E_{\mathrm{hom}}$ any exterior illumination not generated by $\Jimp$, and the
prefactor $-\jj\omega\mu_0$ ties the integral to the locked phasor convention
$\exp(-\jj\omega t)$ \cite{stratton1941}. The primary volume representation is
\begin{equation}
\label{eq:ch2.outgoing-electric-volume-integral}
\E(\mathbf{r})
=
-\jj\omega\mu_0
\int_{\Omega_s}
\overline{\mathbf{G}}_{E,\omega}^{\mathrm{(out)}}(\mathbf{r},\mathbf{r}')\cdot\Jimp(\mathbf{r}')\,d^{3}r'
\;+\;\E_{\mathrm{hom}}(\mathbf{r}),
\end{equation}
\begin{equation}
\label{eq:ch2.outgoing-magnetic-from-curl-integral}
\H(\mathbf{r})
=
\frac{1}{\jj\omega\mu_0}\,\curl\E(\mathbf{r}),
\end{equation}
Equations~\eqref{eq:ch2.outgoing-electric-volume-integral}--\eqref{eq:ch2.outgoing-magnetic-from-curl-integral}
package the dyadic route as an evaluable export map: the electric field is a volume
convolution against the \emph{outgoing} dyad, while the magnetic field is a curl derived
from that electric field in source-free observation regions \cite{stratton1941}. For a
square loop of side \(2a\) with \(k_0 a=1.2\) observed at \(\mathbf{r}\) with
\(k_0|\mathbf{r}|=6\), the integral typically requires \(16\)--\(64\) wavelength-scale
subdivisions of \(\Omega_s\) for \(2\%\) amplitude accuracy when raw
Eq.~\eqref{eq:ch2.free-space-dyadic-green} is used patchwise; self-term corrections from
Subsection~\ref{sec:ch232} dominate error when \(|\mathbf{r}|\to a\)
\cite{harrington2001}. The homogeneous term \(\E_{\mathrm{hom}}\) carries exterior
illumination and must vanish for pure radiation problems \cite{stratton1941}. Equations
\eqref{eq:ch2.outgoing-electric-volume-integral}--\eqref{eq:ch2.outgoing-magnetic-from-curl-integral}
specialize Eqs.~\eqref{eq:ch2.electric-dyadic-convolution}--\eqref{eq:ch2.magnetic-from-electric-dyad}
with outgoing dyads for \(\mathbf{r}\notin\Omega_s\) on simply connected observation
regions; they coincide with Eq.~\eqref{eq:ch2.radiation-integral-outgoing-upgrade} when
\(\E_{\mathrm{hom}}=\mathbf{0}\) \cite{stratton1941}. The prefactor
\(-\jj\omega\mu_0\) matches Eq.~\eqref{eq:ch2.helmholtz-forcing-map} and the locked
phasor axis \cite{stratton1941}. What propagates to a distant probe is the coherent sum of
outgoing dyadic responses $\overline{\mathbf{G}}_{E,\omega}^{\mathrm{(out)}}(\mathbf{r},\mathbf{r}')$
from each source element; what remains reactive at finite $r$ is the non-transverse and
inverse-power content retained in $\mathcal{S}_\omega[\Jimp]$ until $\Pi_{\mathrm{rad}}$ and
regime qualification isolate export \cite{stratton1941}. A near-field scanner
records the full pair $\mathcal{S}_\omega[\Jimp]$; a pattern meter on $\Gamma_R$ with
$k_0 R\gg 1$ approximates $\mathcal{R}_\omega[\Jimp]$ \cite{hansen1988}.
When magnetic impressed sources
\(\mathbf{J}^{\mathrm{(mag)}}_{\mathrm{imp}}\) are primary, duality maps
\(\varepsilon_0\E\leftrightarrow\mu_0\H\) and
\(\Jimp\leftrightarrow\mathbf{J}^{\mathrm{(mag)}}_{\mathrm{imp}}\) yields
\begin{equation}
\label{eq:ch2.outgoing-magnetic-volume-integral}
\H(\mathbf{r})
=
\jj\omega\varepsilon_0
\int_{\Omega_s}
\overline{\mathbf{G}}_{H,\omega}^{\mathrm{(out)}}(\mathbf{r},\mathbf{r}')\cdot
\mathbf{J}^{\mathrm{(mag)}}_{\mathrm{imp}}(\mathbf{r}')\,d^{3}r'
\;+\;\H_{\mathrm{hom}}(\mathbf{r}),
\end{equation}
\begin{equation}
\label{eq:ch2.outgoing-electric-from-curl-magnetic}
\E(\mathbf{r})
=
-\frac{1}{\jj\omega\varepsilon_0}\,\curl\H(\mathbf{r}),
\end{equation}
with \(\overline{\mathbf{G}}_{H,\omega}^{\mathrm{(out)}}\) the magnetic dyad defined
duality-wise from Subsection~\ref{sec:ch231} without re-tabulating its closed form
\cite{stratton1941}. Equations
\eqref{eq:ch2.outgoing-magnetic-volume-integral}--\eqref{eq:ch2.outgoing-electric-from-curl-magnetic}
complete the dual route to Eqs.~\eqref{eq:ch2.outgoing-electric-volume-integral}--\eqref{eq:ch2.outgoing-magnetic-from-curl-integral}
for aperture magnetic currents and dual-equivalent models \cite{stratton1941}.
When both electric and magnetic impressed sources are present, superposition assembles the
full Maxwell pair from the dual volume routes without re-opening gauge or static-potential
paths \cite{stratton1941}. Superposition gives the coupled volume representation
\begin{equation}
\label{eq:ch2.combined-impressed-source-integral}
\begin{aligned}
\E(\mathbf{r})
&=
-\jj\omega\mu_0\!\int_{\Omega_s}\!
\overline{\mathbf{G}}_{E,\omega}^{\mathrm{(out)}}(\mathbf{r},\mathbf{r}')\cdot\Jimp(\mathbf{r}')\,d^{3}r'
-\frac{1}{\jj\omega\varepsilon_0}\,\curl\H^{(\mathrm{mag})}(\mathbf{r})
+\E_{\mathrm{hom}}(\mathbf{r}),\\[4pt]
\H(\mathbf{r})
&=
\jj\omega\varepsilon_0\!\int_{\Omega_s}\!
\overline{\mathbf{G}}_{H,\omega}^{\mathrm{(out)}}(\mathbf{r},\mathbf{r}')\cdot
\mathbf{J}^{\mathrm{(mag)}}_{\mathrm{imp}}(\mathbf{r}')\,d^{3}r'
+\frac{1}{\jj\omega\mu_0}\,\curl\E^{(\mathrm{elec})}(\mathbf{r})
+\H_{\mathrm{hom}}(\mathbf{r}),
\end{aligned}
\end{equation}
where \(\E^{(\mathrm{elec})},\H^{(\mathrm{mag})}\) denote the single-source contributions
from Eqs.~\eqref{eq:ch2.outgoing-electric-volume-integral} and
\eqref{eq:ch2.outgoing-magnetic-volume-integral} \cite{stratton1941}.
Equation~\eqref{eq:ch2.combined-impressed-source-integral} is the first-order Maxwell
integral system on outgoing kernels; it reduces to either branch when one impressed
source vanishes \cite{stratton1941}. Linearity
Eq.~\eqref{eq:ch2.radiation-linearity} follows from Eq.~\eqref{eq:ch2.combined-impressed-source-integral}
on fixed \(\mathcal{B}_\Gamma\) \cite{harrington2001}. The combined integral returns the full
admissible Maxwell pair $\mathcal{S}_\omega[\mathfrak{D}]$: electric and magnetic impressed
sources superpose at observation $\mathbf{r}$ with each contribution evaluated at source
$\mathbf{r}'$ through the outgoing dyads $G$ and $\overline{\mathbf{G}}_{E}$,
$\overline{\mathbf{G}}_{H}$. Pattern, gain, and shell wattmeters do not read this sum directly;
they act on $\mathcal{R}_\omega[\mathfrak{D}]=\Pi_{\mathrm{rad}}\,\mathcal{S}_\omega[\mathfrak{D}]$
after export selection removes reactive and non-propagating fractions that the volume integral
still carries at finite range \cite{stratton1941}.
\begin{figure}[t]
\centering
\ChOneFigBlock{%
\begin{tikzpicture}[ch1fig, node distance=5mm and 7mm]
 \node[ch1box] (jev) {$\Jimp$ volume};
 \node[ch1box, right=7mm of jev] (jmv) {$\mathbf{J}^{\mathrm{(mag)}}_{\mathrm{imp}}$};
 \node[ch1box, right=7mm of jmv] (surf) {Huygens\\$\Gamma_{\mathrm{eq}}$};
 \node[ch1box, below=5mm of jmv] (ge) {$\overline{\mathbf{G}}_{E}^{\mathrm{(out)}}$};
 \node[ch1box, below=5mm of surf] (gh) {$\overline{\mathbf{G}}_{H}^{\mathrm{(out)}}$};
 \node[ch1box, right=7mm of jmv] (ef) {$(\E,\H)$};
 \draw[ch1flow] (jev) -- (ge) -- (ef);
 \draw[ch1flow] (jmv) -- (gh) -- (ef);
 \draw[ch1flow] (surf.east) -- (surf.east |- ef.west);
 \draw[ch1flow] (surf.east |- ef.west) -- (ef.west);
\end{tikzpicture}%
}
\caption{Integral routes: electric volume, magnetic volume, and Huygens surface
representations converge to the same \(\mathcal{S}_\omega\) image on the outgoing branch when
trace matching holds.}
\label{fig:ch2.source-to-field-integral-routes}
\end{figure}
When \(\divg\Jimp-\jj\omega\rho=0\) with nonzero \(\rho\), the complete electric field
includes a scalar-potential gradient driven by \(\rho\):
\begin{equation}
\label{eq:ch2.charge-augmented-electric-integral}
\E(\mathbf{r})
=
-\jj\omega\mu_0\!\int_{\Omega_s}\!
\overline{\mathbf{G}}_{E,\omega}^{\mathrm{(out)}}(\mathbf{r},\mathbf{r}')\cdot\Jimp(\mathbf{r}')\,d^{3}r'
-\grad\!\int_{\Omega_s} G_\omega^{\mathrm{(out)}}(\mathbf{r},\mathbf{r}')\,
\frac{\rho(\mathbf{r}')}{\varepsilon_0}\,d^{3}r'
\;+\;\E_{\mathrm{hom}}(\mathbf{r}),
\end{equation}
with \(G_\omega^{\mathrm{(out)}}\) the outgoing scalar kernel from
Eq.~\eqref{eq:ch2.scalar-green-solution-operator} \cite{stratton1941}. Equation
\eqref{eq:ch2.charge-augmented-electric-integral} enforces Gauss coupling without leaving
\(\mathcal{J}_\Omega\); when Eq.~\eqref{eq:ch1.charge-continuity-harmonic} holds and
\(\rho\) is determined by \(\Jimp\), the gradient term is fixed by the same impressed data
\cite{stratton1941}. The charge-augmented integral still evaluates
$\mathcal{S}_\omega[\mathfrak{D}]$: the scalar kernel $G_\omega^{\mathrm{(out)}}(\mathbf{r},\mathbf{r}')$
propagates charge density from $\mathbf{r}'$ to scalar potential at $\mathbf{r}$, while the
dyadic term propagates $\Jimp$; shell meters on $\mathcal{R}_\omega[\Jimp]$ remain insensitive
to how reactive near storage partitions between the two contributions at finite $r$
\cite{harrington2001}. Lorenz-gauge vector-potential reduction
Eq.~\eqref{eq:ch2.lorenz-field-from-potential} collapses to
Eq.~\eqref{eq:ch2.charge-augmented-electric-integral} on uniform exterior domains; gauge
freedom is not re-opened \cite{harrington2001}.
When the interior of a source region is inaccessible, equivalent currents on a enclosing
surface reproduce the exterior field without resolving volume \(\Jimp\); the surface integrals
below are the operator form of that laboratory substitution \cite{stratton1941}. Let
\(\Gamma_{\mathrm{eq}}\) be a closed surface enclosing \(\Omega_s\), with outward
normal \(\hat{\mathbf{n}}\) and equivalent surface densities
\(\mathbf{J}_{s,\mathrm{eq}}\), \(\mathbf{K}_{s,\mathrm{eq}}\) on
\(\Gamma_{\mathrm{eq}}\) in the trace class of Subsection~\ref{sec:ch211}. The exterior
Huygens representation is
\begin{equation}
\label{eq:ch2.huygens-equivalent-integral}
\E(\mathbf{r})
=
-\jj\omega\mu_0\!\int_{\Gamma_{\mathrm{eq}}}\!
\overline{\mathbf{G}}_{E,\omega}^{\mathrm{(out)}}(\mathbf{r},\mathbf{r}')\cdot
\mathbf{J}_{s,\mathrm{eq}}(\mathbf{r}')\,dS'
+\int_{\Gamma_{\mathrm{eq}}}\!
\nabla'\!G_\omega^{\mathrm{(out)}}(\mathbf{r},\mathbf{r}')\times
\mathbf{K}_{s,\mathrm{eq}}(\mathbf{r}')\,dS',
\end{equation}
\begin{equation}
\label{eq:ch2.huygens-magnetic-integral}
\H(\mathbf{r})
=
\jj\omega\varepsilon_0\!\int_{\Gamma_{\mathrm{eq}}}\!
\overline{\mathbf{G}}_{H,\omega}^{\mathrm{(out)}}(\mathbf{r},\mathbf{r}')\cdot
\mathbf{K}_{s,\mathrm{eq}}(\mathbf{r}')\,dS'
-\int_{\Gamma_{\mathrm{eq}}}\!
\nabla'\!G_\omega^{\mathrm{(out)}}(\mathbf{r},\mathbf{r}')\times
\mathbf{J}_{s,\mathrm{eq}}(\mathbf{r}')\,dS',
\end{equation}
for \(\mathbf{r}\) outside the closed surface bounded by \(\Gamma_{\mathrm{eq}}\)
\cite{stratton1941}. Equations
\eqref{eq:ch2.huygens-equivalent-integral}--\eqref{eq:ch2.huygens-magnetic-integral} are
the surface form of the same outgoing propagation operator;
\(\nabla'\) acts on the source coordinate \(\mathbf{r}'\) \cite{stratton1941}. Tangential
field matching
\begin{equation}
\label{eq:ch2.huygens-tangential-match}
\hat{\mathbf{n}}\times\E\big|_{\Gamma_{\mathrm{eq}}^{+}}
=
\hat{\mathbf{n}}\times\E\big|_{\Gamma_{\mathrm{eq}}^{-}},
\qquad
\hat{\mathbf{n}}\times\H\big|_{\Gamma_{\mathrm{eq}}^{+}}
=
\hat{\mathbf{n}}\times\H\big|_{\Gamma_{\mathrm{eq}}^{-}}
\end{equation}
fixes \(\mathbf{J}_{s,\mathrm{eq}}=\hat{\mathbf{n}}\times(\H^{-}-\H^{+})\),
\(\mathbf{K}_{s,\mathrm{eq}}=\hat{\mathbf{n}}\times(\E^{-}-\E^{+})\) in the stratified
sense of Eq.~\eqref{eq:ch1.interface-conditions} \cite{stratton1941}.
Equivalent-source non-uniqueness Eq.~\eqref{eq:ch2.equivalent-source-ambiguity} applies:
distinct \((\mathbf{J}_{s,\mathrm{eq}},\mathbf{K}_{s,\mathrm{eq}})\) pairs related by
\(\ker\mathcal{R}_\omega\) yield identical far-zone radiation \cite{harrington2001}. The
Huygens surface integrals reproduce exterior fields from equivalent currents on
$\Gamma_{\mathrm{eq}}$: observation $\mathbf{r}$ lies outside the closed surface, source
$\mathbf{r}'$ ranges over $\Gamma_{\mathrm{eq}}$, and the same outgoing kernels
$G_\omega^{\mathrm{(out)}}$ and $\overline{\mathbf{G}}_{E,\omega}^{\mathrm{(out)}}$ propagate
tangential data outward. What is measured on a distant range is still
$\mathcal{R}_\omega[\Jimp]$ after $\Pi_{\mathrm{rad}}$; equivalent currents that differ only by
$\ker\mathcal{R}_\omega$ leave that measurement unchanged while altering reactive near fields
through $\mathcal{S}_\omega$ \cite{stratton1941}.
\phantomsection\label{thm:ch2.source-to-field-integral}
Let \(\mathfrak{D}=(\Jimp,\mathbf{J}^{\mathrm{(mag)}}_{\mathrm{imp}},\rho,\mathcal{B}_\Gamma)\)
be admissible with outgoing radiation condition on exterior problems. Then the physical solution
\(\mathcal{S}_\omega[\mathfrak{D}]=(\E,\H)\) is given by
Eqs.~\eqref{eq:ch2.combined-impressed-source-integral}--\eqref{eq:ch2.charge-augmented-electric-integral}
with outgoing kernels, or by Eqs.~\eqref{eq:ch2.huygens-equivalent-integral}--\eqref{eq:ch2.huygens-tangential-match}
when equivalent surface data replace volume sources \cite{stratton1941}.
The representation agrees with Eq.~\eqref{eq:ch2.outgoing-solution-operator} and
Green--BVP equivalence Eq.~\eqref{eq:ch2.green-bvp-equivalence} on the outgoing branch
fixed by exterior radiative uniqueness (Subsection~\ref{sec:ch243}).
Subsection~\ref{sec:ch231} defined dyadic kernels satisfying
Eq.~\eqref{eq:ch2.dyadic-green-defining-operator}; Subsection~\ref{sec:ch232} justified
convolution on \(\mathcal{J}_\Omega\); Section~\ref{sec:ch24} selected the outgoing
branch. Substituting \(\overline{\mathbf{G}}_{E,\omega}^{\mathrm{(out)}}\) into
Eq.~\eqref{eq:ch2.electric-dyadic-convolution} yields
Eq.~\eqref{eq:ch2.outgoing-electric-volume-integral}. Magnetic and coupled forms follow
from duality and superposition \cite{stratton1941}. Huygens formulas follow by replacing
volume sources with surface distributions that reproduce the same tangential traces
\cite{harrington2001}. Uniqueness on \(\mathcal{F}_{\mathrm{SF}}\) when exterior radiative
uniqueness (Subsection~\ref{sec:ch243}) holds identifies the integral representation
with \(\mathcal{S}_\omega\) \cite{stratton1941}.
The volume and surface integrals above return the full admissible pair $\mathcal{S}_\omega[\mathfrak{D}]$;
pattern meters and gain reports act only after export selection through $\mathcal{R}_\omega$ and
$\Pi_{\mathrm{rad}}$ \cite{stratton1941}. Export selection follows evaluation:
\begin{equation}
\label{eq:ch2.integral-radiation-map}
\mathcal{R}_\omega[\mathfrak{D}]
=
\Pi_{\mathrm{rad}}\,
\Big(\E[\mathfrak{D}],\,\H[\mathfrak{D}]\Big),
\end{equation}
with \(\Pi_{\mathrm{rad}}\) from Eq.~\eqref{eq:ch2.outgoing-radiation-refinement}
\cite{stratton1941}. Pattern cuts, gain reports, and shell power meters act on
Eq.~\eqref{eq:ch2.integral-radiation-map}, not on the raw volume integral alone
\cite{harrington2001}.
On the outgoing branch,
\begin{equation}
\label{eq:ch2.integral-resolvent-equivalence}
\mathcal{S}_\omega[\mathfrak{D}]
=
\mathcal{O}_{\mathrm{out}}\,
\mathcal{L}_\omega^{-1}\!
\begin{bmatrix}\mathbf{s}_\omega(\mathfrak{D})\\ \mathbf{g}_\Gamma\end{bmatrix}
=
\mathcal{O}_{\mathrm{out}}\,
\Big(\mathcal{I}_{\mathrm{vol}}[\mathfrak{D}]\Big),
\end{equation}
with \(\mathcal{I}_{\mathrm{vol}}[\mathfrak{D}]\) the combined volume integrals
Eqs.~\eqref{eq:ch2.combined-impressed-source-integral}--\eqref{eq:ch2.charge-augmented-electric-integral},
connecting integral and operator routes through
Eqs.~\eqref{eq:ch2.resolvent-green-equivalence}--\eqref{eq:ch2.outgoing-resolvent-selection}
\cite{stratton1941}. Equation~\eqref{eq:ch2.integral-resolvent-equivalence}
is the constructive content of Eq.~\eqref{eq:ch2.green-bvp-equivalence} for evaluation
purposes: a method-of-moments matrix implements the resolvent route while patch quadrature
implements the integral route, and both must reference the same outgoing branch
\cite{harrington2001}.
When \(\Omega=\cup_j\Omega_j\) with interface traces from Subsection~\ref{sec:ch234},
Eqs.~\eqref{eq:ch2.combined-impressed-source-integral}--\eqref{eq:ch2.charge-augmented-electric-integral}
hold patchwise with local \(k_j\) and outgoing exterior completion on the infinite exterior
region \cite{stratton1941}. Impressed \(\Jimp\) may vanish in passive exterior regions;
equivalent currents on material boundaries enter through
Eqs.~\eqref{eq:ch2.huygens-equivalent-integral}--\eqref{eq:ch2.huygens-tangential-match}
with \(\Gamma_{\mathrm{eq}}\) on scatterer surfaces \cite{harrington2001}. Interior
cavity problems replace outgoing \(G\) with standing-mode Green functions and require
quotienting Eq.~\eqref{eq:ch2.coercivity-quotient}; those representations are excluded from
the outgoing integral family developed here \cite{stratton1941}.
Equations~\eqref{eq:ch2.outgoing-electric-volume-integral}--\eqref{eq:ch2.charge-augmented-electric-integral}
require linear media, fixed \(\omega\), admissible \(\mathfrak{D}\), outgoing radiation
condition on exterior problems, and observation points outside singular coincidence unless
regularization from Subsection~\ref{sec:ch232} is applied \cite{stratton1941}. Incoming
kernels produce different integrals violating
Eq.~\eqref{eq:ch2.retarded-harmonic-kernel-theorem} \cite{jackson1999}. Anisotropic,
moving, or nonlinear media require modified kernels beyond the locked patchwise-scalar class
\cite{stratton1941}. Near-zone reactive dominance does not invalidate the integral; it
signals that \(\mathcal{R}_\omega[\mathfrak{D}]\) may differ sharply from
\(\mathcal{S}_\omega[\mathfrak{D}]\) until Eq.~\eqref{eq:ch2.regime-operator-screen} is
satisfied \cite{harrington2001}.
One operational distinction closes the display below: volume integrals construct the full Maxwell pair $\mathcal{S}_\omega[\mathfrak{D}]$, whereas pattern meters and watt integrators on $\Gamma_R$ read only after outgoing selection has isolated $\mathcal{R}_\omega[\mathfrak{D}]$. A method-of-moments matrix must not be read as a radiation report; Maxwell theory shows that the matrix implements whichever Green branch and boundary completion were declared, and reactive near-zone content may dominate local probes without contributing to qualified export \cite{stratton1941}. A measurement on a closed scan surface therefore records the complete field state unless a transform and regime qualification explicitly project onto the radiative sector. Synthesis, diagnosis, and inversion campaigns must name the same outgoing branch before any two routes---resolvent inversion and direct quadrature---are compared \cite{harrington2001}. Subsection~\ref{sec:ch252} next fixes distance labels that tell the engineer when far-zone asymptotics may replace the full integral without misreading reactive storage as radiated power.
